%% notes-tex/wet-lab/040-inhibitor-synergy-wetlab/sections/4-results.tex

\section{Results\secstatus{todo}}
\label{sec:results}

\subsection{Claim 1 holds: the combinations grow less than independence predicts and
more than Loewe additivity predicts\secstatus{todo}}
\label{sec:claim1}

\notesec{The single agents} The ex21 titrations give a Hill fit per inhibitor over 27
wells each (Table~\ref{tab:hill}, Figure~\ref{fig:singles}). Potency spans a factor of
50 on a molar basis, from furfural at IC30 9.8 mM to lactic acid at 486 mM, and the
curves are steep: formic acid and lactic acid both sit at the slope bound $h = 20$, as do
most isobole axes, so for those compounds fitness falls from near the top to none between
two adjacent doses. Levulinic acid's fitted IC50 of 10.2 g\,L$^{-1}$ lands on the highest
dose tested, 10 g\,L$^{-1}$, so its curve is extrapolated above that point.

Against the ex23 singles at the ex23 doses, the ex21 fit agrees within the 95 percent
interval for furfural ($+0.03$), acetic acid ($+0.02$), formic acid ($0.00$) and lactic
acid ($-0.02$), and reads high for 5-HMF ($-0.09$, $[-0.12, -0.05]$) and levulinic acid
($-0.09$, $[-0.14, -0.04]$), each over three wells
(\file{results/single\_agent\_ex23\_check.csv}). Acetic acid, dosed in all four runs,
gives IC50 4.00 g\,L$^{-1}$ in ex21 against 2.51, 2.90 and 2.94 in ex26 to ex28, and
ex21's interval $[3.30, 4.37]$ overlaps none of the three isobole intervals. The runs
differ in length and in the ex21 control anomaly, and these data do not identify which
produces the offset.

%% SOURCE: experiments/040-inhibitor-synergy-wetlab/scripts/notes_tex_tables.py ->
%% tables/t4-hill.tex, from results/single_agent_fits.csv (single_agent_curves.py).
\begin{table}[htbp]
\centering
\footnotesize
\caption[]{The ex21 single-agent Hill fits, primary variant (a well that did not grow
enters at fitness 0). The fitted top is the curve's own uninhibited asymptote, used in
place of the ex21 controls. A slope at 20 is at its upper bound and marks a cliff
between two adjacent doses, so the IC50 interval there understates the uncertainty.
Intervals are 95 percent percentiles of 1000 resamples of replicate wells within each
dose.}
\label{tab:hill}
%% GENERATED by experiments/040-inhibitor-synergy-wetlab/scripts/notes_tex_tables.py
%% SOURCE: the result files named in that script's docstring. Do not edit by hand.
\begin{tabular}{lrrrrrlrrl}
\toprule
Inhibitor & Wells & No growth & Fitted top & Hill & IC50 & IC50 95\% & IC30 & IC30 & IC30 95\% \\
 & (n) & (n) & (fitness) & slope $h$ & (g L$^{-1}$) & interval (g L$^{-1}$) & (g L$^{-1}$) & (mM) & interval (mM) \\
\midrule
furfural & 27 & 16 & 1.187 & 2.05 & 1.43 & [1.38, 1.52] & 0.94 & 9.8 & [9.1, 12.2] \\
acetic acid & 27 & 12 & 1.217 & 5.29 & 4.00 & [3.30, 4.37] & 3.41 & 56.7 & [47.3, 64.6] \\
5-HMF & 27 & 12 & 1.002 & 5.16 & 2.66 & [2.59, 2.70] & 2.26 & 17.9 & [17.2, 18.4] \\
formic acid & 27 & 9 & 1.152 & 20.00 (bound) & 2.10 & [2.07, 2.14] & 2.02 & 43.8 & [43.1, 44.6] \\
levulinic acid & 27 & 1 & 1.185 & 5.78 & 10.16 & [8.61, 16.21] & 8.78 & 75.6 & [70.4, 100.0] \\
lactic acid & 27 & 12 & 1.252 & 20.00 (bound) & 45.70 & [44.93, 46.72] & 43.81 & 486.3 & [478.0, 497.1] \\
\bottomrule
\end{tabular}

\end{table}

\tcfig{figures/single_agent_curves.pdf}{%
  \textbf{Each inhibitor's dose response, per run, with the ex23 dose marked.} Wells are
  divided by their own source's fitted top and drawn with that source's Hill curve;
  crosses are wells that did not grow within the run, placed at 0. Stars are the ex23
  single-agent wells, mean and standard deviation of three wells on the ex23 scale where
  the uninhibited control is 1. The dashed vertical line is the ex23 dose, which is where
  every combination in Section~\ref{sec:claim1} sits. The steepness of the acid panels is
  what separates the Bliss and Loewe predictions below.}{fig:singles}
%% Generated by experiments/040-inhibitor-synergy-wetlab/scripts/single_agent_curves.py

\notesec{Potency against the public screen} Only furfural and 5-HMF have a dose in the
corrected Vanacloig store, and the comparison is cross-medium rather than a replication:
Vanacloig dosed the haploid pool at IC30 in anaerobic SYNH3 read at 48 h, the Bioscreen
ran bAID in aerobic YPD to 72 h. Furfural matches within a factor of 1.2, 5-HMF differs
by 5.4 (Table~\ref{tab:vanacloig}). That is the resolution any dose-aware model inherits
from those anchors alone.

%% SOURCE: experiments/040-inhibitor-synergy-wetlab/scripts/notes_tex_tables.py ->
%% tables/t5-vanacloig.tex, from results/vanacloig_ic30.csv (single_agent_curves.py).
\begin{table}[htbp]
\centering
\footnotesize
\caption[]{IC30 in the public screen against IC30 from the ex21 Hill fit. Levulinic acid
has no condition in the corrected store, which the script confirms by reading all 34
condition blocks. Acetate is present only in build 001 and carries no IC30 row here.}
\label{tab:vanacloig}
%% GENERATED by experiments/040-inhibitor-synergy-wetlab/scripts/notes_tex_tables.py
%% SOURCE: the result files named in that script's docstring. Do not edit by hand.
\begin{tabular}{lrrlr}
\toprule
Inhibitor & Vanacloig IC30 & ex21 IC30 & ex21 95\% interval & Ratio ex21 \\
 & (mM) & (mM) & (mM) & / Vanacloig \\
\midrule
furfural & 8.0 & 9.8 & [9.1, 12.2] & 1.23 \\
5-HMF & 3.3 & 17.9 & [17.2, 18.4] & 5.42 \\
levulinic acid & not in store & 75.6 & [70.4, 100.0] &  \\
\bottomrule
\end{tabular}

\end{table}

\notesec{The 63 combinations} Growth is the discriminating readout: 42 of the 63
combinations did not grow within 85 h under the served call, 38 under the software call,
and nothing holding four or more inhibitors grew under either. Loewe additivity
reproduces that pattern and the other two rules do not (Table~\ref{tab:growthscores}).
Loewe reaches accuracy 0.825 and AUROC 0.958 served, and 0.825 and 0.952 under the
software call, with one false positive in both. Bliss from the ex23 singles reaches
0.556 and 0.872 served; Bliss from the ex21 fits and both highest-single-agent variants
predict growth for all 63 combinations, which is accuracy 0.333. Restricted to the 57
combinations of two or more inhibitors, Loewe holds AUROC 0.944 served and 0.939
software against Bliss from the ex23 singles at 0.844 and 0.817.

On the combinations that did grow, the two rules err in opposite directions
(Table~\ref{tab:fitnessscores}, Figure~\ref{fig:obsvspred}). Bliss from the ex23 singles
reaches Spearman 0.573 served and 0.865 software and overestimates fitness by 0.142 and
0.097 on average; Loewe reaches 0.713 and 0.646 and underestimates it by 0.128 and 0.185.
The samples are 21 and 25 grown combinations, and the three-inhibitor rows rest on 3 and
5 of them, too few to read. Loewe's accuracy is also bought partly on extrapolation: for
33 of 63 combinations its solution lies beyond at least one component's tested ex21
range.

%% SOURCE: experiments/040-inhibitor-synergy-wetlab/scripts/notes_tex_tables.py ->
%% tables/t6-growth-scores.tex, from results/mixture_scores.csv (mixture_rules.py).
\begin{table}[htbp]
\centering
\footnotesize
\caption[]{Growth or no growth over all 63 ex23 combinations. A rule predicts growth when
its prediction is at least $\tau$, the lowest fitness of any grown inhibited well under
that call; AUROC uses the prediction as a score and is free of $\tau$. Positives are 21
combinations under the served call and 25 under the software call. Best accuracy and best
AUROC within each call in bold.}
\label{tab:growthscores}
%% GENERATED by experiments/040-inhibitor-synergy-wetlab/scripts/notes_tex_tables.py
%% SOURCE: the result files named in that script's docstring. Do not edit by hand.
\begin{tabular}{llrrrrrrr}
\toprule
Reference rule & Growth & $\tau$ & TP & FP & TN & FN & Accuracy & AUROC \\
 & call & (fitness) & (n) & (n) & (n) & (n) & (fraction) &  \\
\midrule
Bliss, ex21 fits & served & 0.248 & 21 & 42 & 0 & 0 & 0.333 & 0.777 \\
Bliss, ex23 singles & served & 0.248 & 20 & 27 & 15 & 1 & 0.556 & 0.872 \\
Loewe, ex21 fits & served & 0.248 & 11 & 1 & 41 & 10 & \textbf{0.825} & \textbf{0.958} \\
Highest single agent, ex21 fits & served & 0.248 & 21 & 42 & 0 & 0 & 0.333 & 0.663 \\
Highest single agent, ex23 singles & served & 0.248 & 21 & 42 & 0 & 0 & 0.333 & 0.776 \\
Bliss, ex21 fits & software & 0.171 & 25 & 38 & 0 & 0 & 0.397 & 0.741 \\
Bliss, ex23 singles & software & 0.171 & 25 & 38 & 0 & 0 & 0.397 & 0.846 \\
Loewe, ex21 fits & software & 0.171 & 15 & 1 & 37 & 10 & \textbf{0.825} & \textbf{0.952} \\
Highest single agent, ex21 fits & software & 0.171 & 25 & 38 & 0 & 0 & 0.397 & 0.637 \\
Highest single agent, ex23 singles & software & 0.171 & 25 & 38 & 0 & 0 & 0.397 & 0.741 \\
\bottomrule
\end{tabular}

\end{table}

%% SOURCE: experiments/040-inhibitor-synergy-wetlab/scripts/notes_tex_tables.py ->
%% tables/t7-fitness-scores.tex, from results/mixture_scores.csv (mixture_rules.py).
\begin{table}[htbp]
\centering
\footnotesize
\caption[]{Fitness of the combinations that grew, observed as the mean of the grown wells.
The mean of observed minus predicted is signed: Bliss is high on these combinations and
Loewe is low on them. Best Spearman and lowest RMSE within each call in bold. The
software call changes which rule wins the fitness comparison, which is why claim 1 is
reported under both.}
\label{tab:fitnessscores}
%% GENERATED by experiments/040-inhibitor-synergy-wetlab/scripts/notes_tex_tables.py
%% SOURCE: the result files named in that script's docstring. Do not edit by hand.
\begin{tabular}{llrrrr}
\toprule
Reference rule & Growth & Grown & Spearman & RMSE & Mean observed minus \\
 & call & combinations (n) &  & (fitness) & predicted (fitness) \\
\midrule
Bliss, ex21 fits & served & 21 & \textbf{0.717} & 0.259 & -0.167 \\
Bliss, ex23 singles & served & 21 & 0.573 & 0.235 & -0.142 \\
Loewe, ex21 fits & served & 21 & 0.713 & \textbf{0.208} & +0.128 \\
Highest single agent, ex21 fits & served & 21 & 0.633 & 0.267 & -0.183 \\
Highest single agent, ex23 singles & served & 21 & 0.476 & 0.243 & -0.162 \\
Bliss, ex21 fits & software & 25 & 0.525 & 0.238 & -0.118 \\
Bliss, ex23 singles & software & 25 & \textbf{0.865} & \textbf{0.137} & -0.097 \\
Loewe, ex21 fits & software & 25 & 0.646 & 0.279 & +0.185 \\
Highest single agent, ex21 fits & software & 25 & 0.431 & 0.252 & -0.149 \\
Highest single agent, ex23 singles & software & 25 & 0.824 & 0.198 & -0.158 \\
\bottomrule
\end{tabular}

\end{table}

\tcfig{figures/mixture_rules_observed_vs_predicted.pdf}{%
  \textbf{Observed against predicted fitness for the 63 combinations, by rule and growth
  call.} Columns are the reference rules of Section~\ref{sec:rules} and rows are the two
  growth calls. Crosses are combinations with no growth within 85 h, drawn at 0, and the
  dashed lines mark that call's $\tau$. The Bliss and highest-single-agent columns place
  nearly every combination above $\tau$, which is the accuracy 0.333 of
  Table~\ref{tab:growthscores} seen directly; the Loewe column puts the no-growth set
  below $\tau$ and the grown set above its own prediction.}{fig:obsvspred}
%% Generated by experiments/040-inhibitor-synergy-wetlab/scripts/mixture_rules.py

\notesec{The pairs} Twelve of the 15 pairs deviate below Bliss with a bootstrap interval
excluding zero, two are additive and one is antagonistic
(Table~\ref{tab:pairs}, Figure~\ref{fig:pairdev}). The largest deviation is formic acid
with lactic acid, which grew in no well under either call ($-0.978$ served). The two
additive pairs are 5-HMF with acetic acid and 5-HMF with formic acid, under both calls.
The one antagonistic pair, 5-HMF with furfural at $+0.058$ $[+0.018, +0.103]$ served,
flips to synergy under the software call at $-0.040$ $[-0.052, -0.026]$, and the two
5-HMF pairs whose wells the served call reads as no growth shrink by a factor of 3 under
the software call. Every 5-HMF statement in this paragraph is therefore call dependent,
and that is the same 12 wells of Section~\ref{sec:data}.

%% SOURCE: experiments/040-inhibitor-synergy-wetlab/scripts/notes_tex_tables.py ->
%% tables/t8-pairs.tex, from results/pair_deviation.csv (mixture_rules.py).
\begin{table}[htbp]
\centering
\scriptsize
\caption[]{The 15 pairs against Bliss independence from the observed ex23 singles, ordered
by the served deviation. Observed is the mean of three wells with a well that did not grow
counted as 0. The interval is 95 percent from 2000 resamples of the replicate wells of the
pair and of both singles, and is shown for the served call; the call is synergy when the
interval lies below zero, antagonism when above, additive otherwise. The software
deviations carry their own intervals in \file{results/pair\_deviation.csv}.}
\label{tab:pairs}
%% GENERATED by experiments/040-inhibitor-synergy-wetlab/scripts/notes_tex_tables.py
%% SOURCE: the result files named in that script's docstring. Do not edit by hand.
\begin{tabular}{llrrrlcrc}
\toprule
Pair & Doses & Observed & Bliss & Deviation & 95\% & Call, & Deviation, & Call, \\
 & (g L$^{-1}$) & (fitness) & (fitness) & served & interval & served & software & software \\
\midrule
formic acid + lactic acid & 1.00 / 20.00 & 0.000 & 0.978 & -0.978 & [-1.023, -0.934] & synergy & -0.718 & synergy \\
acetic acid + formic acid & 2.00 / 1.00 & 0.443 & 0.994 & -0.551 & [-0.596, -0.500] & synergy & -0.286 & synergy \\
5-HMF + lactic acid & 2.52 / 20.00 & 0.000 & 0.471 & -0.471 & [-0.497, -0.447] & synergy & -0.147 & synergy \\
acetic acid + levulinic acid & 2.00 / 6.00 & 0.399 & 0.867 & -0.468 & [-0.613, -0.358] & synergy & -0.247 & synergy \\
5-HMF + levulinic acid & 2.52 / 6.00 & 0.000 & 0.417 & -0.417 & [-0.434, -0.400] & synergy & -0.152 & synergy \\
formic acid + levulinic acid & 1.00 / 6.00 & 0.455 & 0.865 & -0.411 & [-0.555, -0.312] & synergy & -0.183 & synergy \\
acetic acid + lactic acid & 2.00 / 20.00 & 0.622 & 0.980 & -0.359 & [-0.475, -0.234] & synergy & -0.201 & synergy \\
furfural + lactic acid & 1.50 / 20.00 & 0.292 & 0.498 & -0.206 & [-0.254, -0.155] & synergy & -0.124 & synergy \\
lactic acid + levulinic acid & 20.00 / 6.00 & 0.692 & 0.854 & -0.162 & [-0.231, -0.097] & synergy & -0.101 & synergy \\
furfural + levulinic acid & 1.50 / 6.00 & 0.281 & 0.440 & -0.160 & [-0.204, -0.113] & synergy & -0.146 & synergy \\
acetic acid + furfural & 2.00 / 1.50 & 0.417 & 0.506 & -0.089 & [-0.139, -0.038] & synergy & -0.188 & synergy \\
formic acid + furfural & 1.00 / 1.50 & 0.434 & 0.505 & -0.071 & [-0.122, -0.017] & synergy & -0.219 & synergy \\
5-HMF + acetic acid & 2.52 / 2.00 & 0.449 & 0.479 & -0.029 & [-0.081, 0.015] & additive & +0.015 & additive \\
5-HMF + formic acid & 2.52 / 1.00 & 0.461 & 0.478 & -0.017 & [-0.057, 0.023] & additive & +0.014 & additive \\
5-HMF + furfural & 2.52 / 1.50 & 0.301 & 0.243 & +0.058 & [0.018, 0.103] & antagonism & -0.040 & synergy \\
\bottomrule
\end{tabular}

\end{table}

\tcfig{figures/mixture_rules_pair_deviation.pdf}{%
  \textbf{Observed minus Bliss for the 15 pairs, under both growth calls.} Bars are 95
  percent bootstrap intervals over resampled replicate wells. Orange takes Bliss from the
  observed ex23 singles, purple from the ex21 Hill fits with the fit-draw uncertainty
  included, which is why the purple intervals are wider. A bar whose interval clears zero
  is a synergy or antagonism call in Table~\ref{tab:pairs}; the 5-HMF pairs are the ones
  that move between the two calls.}{fig:pairdev}
%% Generated by experiments/040-inhibitor-synergy-wetlab/scripts/mixture_rules.py

\notesec{The isoboles} The three dose grids put the same question on a surface rather
than at a point (Table~\ref{tab:isoboles}, Figure~\ref{fig:isoboles}). Furfural with
acetic acid lies below Bliss by 0.082 $[-0.142, -0.017]$ over its 81 interior cells, and
formic acid with acetic acid by 0.446 $[-0.536, -0.367]$, with 31 interior cells where
Bliss predicts growth and nothing grew. The 5-HMF grid is additive against Bliss,
$-0.017$ $[-0.043, +0.007]$, and is flagged rather than trusted: one run, with growth
islands at 5-HMF 2.0 g\,L$^{-1}$. No interior cell of any grid grew less than Loewe
predicts, while 10, 3 and 12 cells grew where Loewe predicts none, so the observed front
sits between the references and closest to Loewe, the same ordering the 63 combinations
give.

%% SOURCE: experiments/040-inhibitor-synergy-wetlab/scripts/notes_tex_tables.py ->
%% tables/t9-isoboles.tex, from results/isobole_summary.csv (mixture_rules.py).
\begin{table}[htbp]
\centering
\scriptsize
\caption[]{The three isobole grids, served call. Excess over Bliss is computed against the
empirical Bliss surface built from each grid's own axes, averaged over the 81 interior
cells, with a 95 percent interval from 2000 resamples of the two plate wells per cell.
The last three columns count interior cells that contradict a reference in one direction
only. ex28 is analyzed and flagged, not trusted.}
\label{tab:isoboles}
%% GENERATED by experiments/040-inhibitor-synergy-wetlab/scripts/notes_tex_tables.py
%% SOURCE: the result files named in that script's docstring. Do not edit by hand.
\begin{tabular}{lrrrlcrrr}
\toprule
Run and pair & Interior & Grew & Mean excess & 95\% & Call against & Bliss grows, & Loewe grows, & Loewe none, \\
 & cells (n) & (n) & over Bliss & interval & Bliss & none grew & none grew & growth seen \\
 &  &  & (fitness) &  &  & (n) & (n) & (n) \\
\midrule
ex26, furfural x acetic acid & 81 & 14 & -0.082 & [-0.142, -0.017] & synergy & 1 & 0 & 10 \\
ex27, formic acid x acetic acid & 81 & 24 & -0.446 & [-0.536, -0.367] & synergy & 31 & 0 & 3 \\
ex28, 5-HMF x acetic acid (flagged) & 81 & 17 & -0.017 & [-0.043, 0.007] & additive & 3 & 0 & 12 \\
\bottomrule
\end{tabular}

\end{table}

\tcfig{figures/mixture_rules_isoboles.pdf}{%
  \textbf{The three isobole grids, observed and against Bliss.} Top row: observed fitness
  per dose cell, mean of two plates with a cell that did not grow at 0, the black contour
  around the cells that grew and the dashed blue line the Loewe additive front at that
  run's $\tau$ from the grid-axis Hill fits. Bottom row: observed minus the empirical
  Bliss surface, so blue is growth below independence. The formic acid grid is where the
  two references separate most: its observed front follows the Loewe line while 31
  interior cells that Bliss expects to grow did not.}{fig:isoboles}
%% Generated by experiments/040-inhibitor-synergy-wetlab/scripts/mixture_rules.py

Claim 1 is therefore supported in the direction registered, synergy relative to Bliss,
with one correction to the registration: the number of grown combinations is 21 under the
primary served call rather than the 25 the claim names, which is the software call's
count. Hypothesis (untested): the steep single-agent curves, four of six at or near the
slope bound, are what drive the Bliss and Loewe predictions this far apart. No
calculation here separates that from the alternatives.

\subsection{Claim 2 fails as registered: profile similarity does not anticipate which
pairs deviate\secstatus{todo}}
\label{sec:claim2}

The registered expectation was that two inhibitors with similar deletion profiles, and so
with shared targets, would be additive, and dissimilar ones would deviate. With the
measured profiles in the matrix nothing in Table~\ref{tab:similarity} supports it: the
largest association is $+0.311$ at $p = 0.26$, every permutation $p$ is above 0.26, and
the sign flips between the centered Spearman and the top-100 Jaccard measure. Only 6 of
the 15 pairs have two predicted profiles in that matrix, so the test is thin as well.

The only significant association is in the all-predicted matrix and has the opposite sign
to the registration: the more similar two predicted profiles, the more synergistic the
pair (raw Spearman $-0.686$, $p = 0.006$ under the software call; $-0.493$, $p = 0.066$
served). It is a chemistry-class regularity rather than a chemogenomic one. The three
small acids have near-duplicate predicted profiles, centered Spearman 0.91 to 0.99, which
is fingerprint similarity carried through a ridge fit whose 32 compounds include no C1 to
C3 acid, and they are also the strongly synergistic pairs: formic with lactic acid grew in
no well, acetic with formic acid deviates by $-0.551$, acetic with levulinic acid by
$-0.468$. The association therefore says that weak-acid pairs synergize, which the
fingerprint already encodes, and it depends on the growth call ($p = 0.07$ to $p = 0.43$
served). The three isoboles cannot separate the two readings: furfural with acetic acid
has similarity 0.04 in the best-available matrix and 0.51 in the all-predicted one and is
mildly synergistic, while formic with acetic acid has $-0.08$ and 0.99 and is strongly
synergistic.

%% SOURCE: experiments/040-inhibitor-synergy-wetlab/scripts/notes_tex_tables.py ->
%% tables/t10-similarity.tex, from results/similarity_vs_deviation.csv
%% (similarity_vs_deviation.py).
\begin{table}[htbp]
\centering
\footnotesize
\caption[]{Profile similarity against the pair deviation from Bliss, over the 15 pairs. A
negative Spearman means similar profiles deviate further below Bliss, which is the
opposite of the registered direction. Permutation $p$ is over 10,000 draws and is set in
bold below 0.05. The last column counts the pairs whose two profiles are both
predictions, which is where the fingerprint regularity of the small acids enters.}
\label{tab:similarity}
%% GENERATED by experiments/040-inhibitor-synergy-wetlab/scripts/notes_tex_tables.py
%% SOURCE: the result files named in that script's docstring. Do not edit by hand.
\begin{tabular}{lllrrrlr}
\toprule
Profile & Similarity & Growth & Pairs & Spearman & Permutation & Leave-one-pair-out & Both profiles \\
matrix & measure & call & (n) & against deviation & $p$ & range (Spearman) & predicted (n) \\
\midrule
best available & centered Spearman & served & 15 & +0.204 & 0.462 & [+0.081, +0.481] & 6 \\
best available & centered Spearman & software & 15 & +0.311 & 0.263 & [+0.226, +0.613] & 6 \\
best available & raw Spearman & served & 15 & -0.179 & 0.522 & [-0.319, +0.011] & 6 \\
best available & raw Spearman & software & 15 & +0.086 & 0.768 & [+0.007, +0.336] & 6 \\
best available & top-100 Jaccard, centered & served & 15 & -0.269 & 0.328 & [-0.426, -0.099] & 6 \\
best available & top-100 Jaccard, centered & software & 15 & -0.022 & 0.940 & [-0.113, +0.205] & 6 \\
all predicted & centered Spearman & served & 15 & -0.225 & 0.425 & [-0.354, -0.051] & 15 \\
all predicted & centered Spearman & software & 15 & -0.525 & 0.050 & [-0.662, -0.420] & 15 \\
all predicted & raw Spearman & served & 15 & -0.493 & 0.066 & [-0.653, -0.380] & 15 \\
all predicted & raw Spearman & software & 15 & -0.686 & \textbf{0.006} & [-0.771, -0.618] & 15 \\
all predicted & top-100 Jaccard, centered & served & 15 & -0.318 & 0.246 & [-0.462, -0.166] & 15 \\
all predicted & top-100 Jaccard, centered & software & 15 & -0.709 & \textbf{0.005} & [-0.780, -0.648] & 15 \\
\bottomrule
\end{tabular}

\end{table}

\tcfig{figures/similarity_vs_deviation.pdf}{%
  \textbf{Pair deviation from Bliss against profile similarity, by matrix and growth
  call.} One point per pair. The best-available panels carry the measured furfural and
  Hoepfner acetate columns and show no association; the all-predicted panels show the
  negative trend of Table~\ref{tab:similarity}, which the small-acid pairs at high
  similarity and large negative deviation produce.}{fig:simdev}
%% Generated by experiments/040-inhibitor-synergy-wetlab/scripts/similarity_vs_deviation.py

\tcfig{figures/profile_similarity_centered.pdf}{%
  \textbf{Centered profile similarity among the six inhibitors, in both matrices.}
  Spearman correlation on a diverging scale and top-100 Jaccard overlap on a sequential
  scale, over the genes each pair of profiles shares. In the all-predicted matrix the
  three small acids form a near-duplicate block at 0.91 to 0.99. In the best-available
  matrix the Hoepfner acetate column is near zero against everything, at most 0.09 in
  absolute value, which is the measured cross-screen column failing to resemble any
  prediction.}{fig:profilesim}
%% Generated by experiments/040-inhibitor-synergy-wetlab/scripts/inhibitor_profiles.py

\subsection{Claim 3 holds on public data: deletion profiles compose as the mean of the
singles\secstatus{todo}}
\label{sec:claim3}

Over Hillenmeyer's 26 two-compound environments the mean of the two matched single
profiles has the highest $R^2$ of the three fixed rules in 24 of 26 pairs, and in 20 of
the 22 that are not flagged (Table~\ref{tab:hetrules}, Figure~\ref{fig:hetrulefit}). The
sum wins once and the maximum once, in a pair where every rule is below 0.14. Medians over
the 26 are $R^2$ 0.365 for the mean, 0.091 for the sum and 0.050 for the maximum, against
0.458 for the per-pair linear fit, whose median coefficients are $a = 0.529$ and
$b = 0.593$ with $a + b = 1.06$ and an intercept near zero. The finding is about magnitude
and not about ordering: the sum and the mean order genes identically and share one
Spearman of 0.579, so the sum's low $R^2$ is entirely its overshoot.

The matched singles do carry the pair. That Spearman of 0.579 compares with 0.066 for two
random YPD singles and 0.114 for singles matched on the control set, and exceeds the
matched null in 24 of 26 pairs at $p \le 0.05$ over 200 draws
(Table~\ref{tab:hetemergent}); both exceptions are pairs whose azole single sits a
twofold step from the pair dose.

%% SOURCE: experiments/040-inhibitor-synergy-wetlab/scripts/notes_tex_tables.py ->
%% tables/t11-het-rules.tex, from results/het_rule_fit.csv and results/het_summary.json
%% (het_mixture_composition.py).
\begin{table}[htbp]
\centering
\footnotesize
\caption[]{How a two-compound profile relates to its two single profiles, over
Hillenmeyer's 26 two-compound environments and the 22 that are not flagged. $R^2$ is of
the rule as it stands, with no refit, so a rule of the wrong magnitude scores low or
negative; Spearman is scale-free, which is why the sum and the mean share one value. The
sum is the Bliss product on this scale (Equation~\ref{eq:blisslog}). Best value per
column in bold. The free linear fit is fitted per pair and so is an upper reference, not a
rule that could be applied to an unmeasured pair.}
\label{tab:hetrules}
%% GENERATED by experiments/040-inhibitor-synergy-wetlab/scripts/notes_tex_tables.py
%% SOURCE: the result files named in that script's docstring. Do not edit by hand.
\begin{tabular}{lrrrrl}
\toprule
Composition rule & Median $R^2$ & Median $R^2$ & Median Spearman & Median Spearman & Highest $R^2$ \\
 & 26 pairs & 22 unflagged & 26 pairs & 22 unflagged & in \\
\midrule
Sum (Bliss) & 0.091 & 0.119 & 0.579 & 0.587 & 1 of 26 \\
Mean & 0.365 & 0.417 & 0.579 & 0.587 & 24 of 26 \\
Max & 0.050 & 0.106 & 0.411 & 0.482 & 1 of 26 \\
Free linear & \textbf{0.458} & \textbf{0.469} & \textbf{0.634} & \textbf{0.595} & fitted per pair \\
\bottomrule
\end{tabular}
%% median free-linear coefficients over 26 pairs: a 0.529, b 0.593, a+b 1.056, c 0.014

\end{table}

\tcfig{figures/het_rule_fit.pdf}{%
  \textbf{(a)} $R^2$ of each composition rule against the observed pair profile, one
  group of bars per two-compound environment, grouped into pair families by the vertical
  rules; an $R^2$ below $-1$ is drawn as a down triangle at the floor. The mean is highest
  in all but two groups. \textbf{(b)} The same comparison by Spearman, where the sum and
  the mean coincide, with open circles giving the Spearman between the two matched single
  profiles, median 0.254. Gray bars are the four flagged tacrolimus pairs, whose nearest
  dosed single is 402 to 804-fold away. 5,677 to 5,704 genes per
  pair.}{fig:hetrulefit}
%% Generated by experiments/040-inhibitor-synergy-wetlab/scripts/het_mixture_composition.py

\notesec{Emergent and masked genes sit inside the noise} Pooled over the 26 pairs, 6,234
of 16,172 pair hits are emergent (0.385) and 18,341 of 28,279 single-union hits are
masked (0.649), and the median emergent fraction of 0.355 is well below both nulls, 0.621
random and 0.613 control-set matched, in 24 of 26 pairs. The rate cannot be read as
biology, because the only measurable noise floor is of the same size: the two
near-replicate single environments in the release agree at Spearman 0.242 and 0.103, and
54 and 64 percent of one member's hits are not hits in the other
(Table~\ref{tab:hetemergent}). With a $|z| > 2$ cut on single arrays, emergent and masked
counts are dominated by irreproducibility, and these data do not separate the two. The
floor itself carries two caveats: $n = 2$, and the near-replicates span different control
sets while every pair shares its control set with its singles, so the like-for-like floor
is not measured.

Both fractions fall with the similarity of the two singles, Spearman $-0.657$
($p = 2.6 \times 10^{-4}$) for emergent and $-0.786$ ($p = 2.0 \times 10^{-6}$) for
masked over 26 pairs (Figure~\ref{fig:hetemergent}). Hypothesis (untested): two similar
singles share hits, so their union covers more of the pair's hits and leaves fewer to be
masked, which would produce this trend with no interaction at all.

%% SOURCE: experiments/040-inhibitor-synergy-wetlab/scripts/notes_tex_tables.py ->
%% tables/t12-het-emergent.tex, from results/het_emergent_masked.csv,
%% results/het_near_replicates.csv and results/het_summary.json
%% (het_mixture_composition.py).
\begin{table}[htbp]
\centering
\scriptsize
\caption[]{Hit-level composition against the two nulls, and the noise floor that bounds
what it can mean. A hit is robust $|z| > 2$ within an environment. The nulls replace the
matched singles by random YPD singles, 200 draws, the second restricted to singles
sharing a control set with the pair. The last two rows are the only near-replicate single
environments in the release, two 5-fluorouracil doses within 1 percent; their hit
discordance exceeds the median emergent fraction above.}
\label{tab:hetemergent}
%% GENERATED by experiments/040-inhibitor-synergy-wetlab/scripts/notes_tex_tables.py
%% SOURCE: the result files named in that script's docstring. Do not edit by hand.
\begin{tabular}{p{44mm}rrrlr}
\toprule
Quantity & Median & Median, & Median, & Pairs beyond & Pooled, \\
 & observed & random & matched & matched null & 26 pairs \\
 &  & singles & singles & ($p \le 0.05$) &  \\
\midrule
Emergent fraction (of pair hits) & 0.355 & 0.621 & 0.613 & 24 of 26 & 0.385 \\
Masked fraction (of single-union hits) & 0.623 &  &  &  & 0.649 \\
Sum-rule Spearman against the pair profile & 0.579 & 0.066 & 0.114 & 24 of 26 &  \\
Near-replicate singles, 38.4 against 38.5 $\mu$M 5-fluorouracil & 0.242 &  &  &  & hits not shared 0.603 / 0.536 \\
Near-replicate singles, 76.9 against 76.8 $\mu$M 5-fluorouracil & 0.103 &  &  &  & hits not shared 0.483 / 0.640 \\
\bottomrule
\end{tabular}

\end{table}

\tcfig{figures/het_emergent_masked.pdf}{%
  \textbf{(a)} Emergent fraction against the similarity of the two single profiles, one
  point per two-compound environment, with the two null means drawn as dashes and crosses
  at each pair's own similarity. Every observed point lies below both nulls.
  \textbf{(b)} Masked fraction against the same similarity, falling more steeply,
  Spearman $-0.786$. Open symbols are the four flagged tacrolimus
  pairs.}{fig:hetemergent}
%% Generated by experiments/040-inhibitor-synergy-wetlab/scripts/het_mixture_composition.py

\notesec{The weights move with dose} On the one dose grid HET provides, methotrexate
against 5-fluorouracil at three doses each, the mean is the best fixed rule at all nine
dose pairs, $R^2$ 0.223 to 0.526, while the sum ranges from $-0.724$ to 0.305
(Figure~\ref{fig:hetgrid}). The free-linear coefficients are not stable across the grid:
$a$ runs 0.128 to 0.757 and $b$ 0.239 to 0.765, with the methotrexate weight high at
methotrexate 500 $\mu$M and the 5-fluorouracil weight dominating at 125 $\mu$M. So the
mean is a default rather than a constant. Hypothesis (untested): the weight follows
whichever partner is relatively stronger at that dose ratio; no single-agent potency was
computed on these environments to test it.

\tcfig{figures/het_mtx_5fu_grid.pdf}{%
  \textbf{The methotrexate by 5-fluorouracil dose grid, nine dose pairs.}
  \textbf{(a)} The free-linear weight $a$ on the methotrexate single, which rises with
  the methotrexate dose. \textbf{(b)} The weight $b$ on the 5-fluorouracil single, which
  dominates at the lower methotrexate dose. \textbf{(c)} $R^2$ of the free linear fit,
  0.247 to 0.559, lowest at methotrexate 500 $\mu$M. \textbf{(d)} $R^2$ of the fixed mean
  rule, 0.223 to 0.526, which is below the free fit at every cell and above the sum at
  every cell.}{fig:hetgrid}
%% Generated by experiments/040-inhibitor-synergy-wetlab/scripts/het_mixture_composition.py

Carrying the mean rule to the six inhibitors is a labeled hypothesis (untested). HET is a
heterozygous diploid pool in aerobic YPD read as $\log_2$ depletion over 20 generations,
the inhibitor profiles are haploid homozygous deletions in anaerobic SYNH3 at IC30, and
no two-compound inhibitor screen exists to check it. The rule is also a statement about
the shape of a deletion profile, not about the host's growth, which is what
Section~\ref{sec:claim1} scores.

\subsection{Claim 4, post-analysis: the predicted profiles carry only the generic stress
component\secstatus{todo}}
\label{sec:claim4}

\notesec{What the predictions are worth} Leave-one-compound-out over the 32 corrected
compounds gives a median centered Spearman of 0.387, interquartile range 0.241 to 0.482,
which is above experiment 038's bar of 0.359 with 31 rather than 25 to 26 compounds
fitted, and 038's own fold means for these two compounds, 0.357 and 0.392, reproduce
within that spread (Table~\ref{tab:profiles}, Figure~\ref{fig:profpred}). The two
compound-level numbers are not equally meaningful. Furfural reaches 0.341 against a
profile whose replicate reliability is 0.769. 5-HMF reaches 0.425 against a profile whose
replicate reliability is $-0.048$, so that score is agreement with noise plus the shared
stress component and is not evidence that 5-HMF biology is predicted.

Acetate shows the limit from the other side. The predicted acetic acid profile correlates
with Hoepfner's measured sodium acetate profile at $-0.084$ over 3,453 genes, and the two
measured acetate profiles, Vanacloig build 001 and Hoepfner, do not agree with each other
either, $-0.011$ centered. The acetic acid and sodium acetate predictions are nearly the
same profile, 0.932 centered, so the salt-for-acid substitution of
Section~\ref{sec:data} changes little on the prediction side and the disagreement is
between screens; which screen bAID in YPD resembles is untested. Build 001 and build 002
give the same furfural profile, raw Spearman 0.99998, so the normalization change between
builds separates nothing here.

%% SOURCE: experiments/040-inhibitor-synergy-wetlab/scripts/notes_tex_tables.py ->
%% tables/t13-profiles.tex, from results/profile_predicted_vs_measured.csv
%% (inhibitor_profiles.py).
\begin{table}[htbp]
\centering
\footnotesize
\caption[]{Compound-cold profile checks. Centering subtracts each gene's mean over the 32
build-002 compounds, which removes the component every compound shares, so the centered
column is the compound-specific agreement and the raw column is inflated by that shared
part. A score against a profile whose replicate reliability is near or below zero is a
score against noise.}
\label{tab:profiles}
%% GENERATED by experiments/040-inhibitor-synergy-wetlab/scripts/notes_tex_tables.py
%% SOURCE: the result files named in that script's docstring. Do not edit by hand.
\begin{tabular}{p{66mm}lrlrr}
\toprule
Comparison & Genes & Centered & IQR & Raw & Measured replicate \\
 & (n) & Spearman & (Spearman) & Spearman & reliability (Spearman) \\
\midrule
Leave-one-compound-out over 32 compounds, median & 32 compounds & 0.387 & [0.241, 0.482] &  &  \\
furfural, leave-one-out against build 002 & 3492 & 0.341 &  & 0.258 & 0.769 \\
5-HMF, leave-one-out against build 002 & 3503 & 0.425 &  & 0.260 & -0.048 \\
acetic acid predicted against Hoepfner sodium acetate & 3453 & -0.084 &  & 0.043 &  \\
acetic acid prediction against sodium acetate prediction & 3587 & 0.932 &  & 0.990 &  \\
levulinic acid predicted against build 001 & 3485 & 0.190 &  & 0.255 & 0.147 \\
Vanacloig build 001 sodium acetate against Hoepfner sodium acetate & 3376 & -0.011 &  & 0.034 &  \\
furfural, build 001 against build 002 & 3492 & 0.998 &  & 1.000 & 0.769 \\
\bottomrule
\end{tabular}

\end{table}

\tcfig{figures/profile_predicted_vs_measured.pdf}{%
  \textbf{Predicted against measured centered profiles, one point per deletion strain.}
  Panels are furfural and 5-HMF with themselves held out of the fit, acetic acid against
  the Hoepfner sodium acetate profile, and the supplementary build-001 levulinic acid
  row. The furfural panel is the one comparison whose measured axis has a replicate
  reliability above 0.5; the acetate panel is the cross-screen comparison that fails,
  and its scatter is round rather than elongated.}{fig:profpred}
%% Generated by experiments/040-inhibitor-synergy-wetlab/scripts/inhibitor_profiles.py

\notesec{The enrichment test, against the control that matters} A nominated set is
enriched for something in every case, so the test is whether it is enriched for anything
the 32-compound gene mean does not already give (Table~\ref{tab:go},
Figure~\ref{fig:go}). It is not. The measured furfural profile's top 100 genes carry no
term at FDR 0.05, the best being ribosome at 0.057. The predicted profiles carry terms,
but control terms: 19 of 20 for furfural, 40 of 43 for levulinic acid, 48 of 49 for
acetic acid. The 11 terms the acetic acid prediction shares with the Hoepfner profile are
all in the control too, vesicle transport and vesicle tethering complex among them. On
centered profiles no compound shares a single significant term between measured and
predicted.

%% SOURCE: experiments/040-inhibitor-synergy-wetlab/scripts/notes_tex_tables.py ->
%% tables/t14-go.tex, from results/go_overlap_measured_vs_predicted.csv
%% (inhibitor_nominations.py).
\begin{table}[htbp]
\centering
\scriptsize
\caption[]{GO terms at FDR 0.05 over the top 100 most sensitive genes, for each compound
with a measured profile and for its compound-cold prediction. The control is the
build-002 gene mean over the 32 compounds, which is what a model with no
compound-specific signal outputs; the control columns are computed on raw profiles only.
A term in the predicted-and-control column is not evidence about the compound.}
\label{tab:go}
%% GENERATED by experiments/040-inhibitor-synergy-wetlab/scripts/notes_tex_tables.py
%% SOURCE: the result files named in that script's docstring. Do not edit by hand.
\begin{tabular}{llrrrrrrr}
\toprule
Measured profile & Target & Shared & Terms, & Terms, & Shared & Terms, & Measured and & Predicted and \\
 &  & top genes & measured & predicted & terms & control & control & control \\
 &  & (n) & (n) & (n) & (n) & (n) & (n) & (n) \\
\midrule
furfural, build 002 & raw & 19 & 0 & 20 & 0 & 58 & 0 & 19 \\
5-HMF, build 002 (noise) & raw & 22 & 0 & 4 & 0 & 58 & 0 & 4 \\
acetic acid, Hoepfner & raw & 10 & 20 & 49 & 11 & 58 & 11 & 48 \\
levulinic acid, build 001 & raw & 24 & 0 & 43 & 0 & 58 & 0 & 40 \\
furfural, build 002 & centered & 15 & 0 & 13 & 0 &  &  &  \\
5-HMF, build 002 (noise) & centered & 43 & 15 & 6 & 0 &  &  &  \\
acetic acid, Hoepfner & centered & 15 & 5 & 5 & 0 &  &  &  \\
levulinic acid, build 001 & centered & 38 & 11 & 0 & 0 &  &  &  \\
\bottomrule
\end{tabular}

\end{table}

\tcfig{figures/go_measured_vs_predicted_centered.pdf}{%
  \textbf{The eight strongest terms of each measured centered profile, beside the
  prediction and the control.} Bars are $-\log_{10}$ of the Benjamini-Hochberg FDR for
  the measured profile, for its compound-cold prediction, and for the 32-compound gene
  mean; a zero bar means the term was neither significant nor among that set's ten best.
  No compound has a term where the measured and predicted bars are both
  significant.}{fig:go}
%% Generated by experiments/040-inhibitor-synergy-wetlab/scripts/inhibitor_nominations.py

\notesec{The nominations, which are hypotheses} Composing the six single columns by the
sum rule and categorizing gives 60 to 68 core genes per combination in the
best-available matrix and 105 in the all-predicted one (Table~\ref{tab:nominations}). The
core set is led by Golgi-transport and chromatin genes, \gene{YPT6}, \gene{RIC1},
\gene{VPS63}, \gene{COG5}, \gene{GOS1}, \gene{SNF6} and \gene{DEG1}, which are the genes
sick under most Vanacloig compounds. Conflict genes, sensitive in one member and
resistant in another, number 5 to 25 per combination in the best-available matrix and
exactly 1 in the all-predicted one, because predicted profiles are too alike to disagree.
No screen of the bAID pool under any combination exists, so every one of these is a
hypothesis for an experiment that has not been run, and a nominated set's enrichment is
not evidence of compound biology unless it exceeds the control above.

%% SOURCE: experiments/040-inhibitor-synergy-wetlab/scripts/notes_tex_tables.py ->
%% tables/t15-nominations.tex, from results/nominations_summary.csv
%% (inhibitor_nominations.py).
\begin{table}[htbp]
\centering
\footnotesize
\caption[]{Nominated gene counts over the 35 combinations of two or three inhibitors, rule
sum. The core set is defined on the six single columns and so varies only with which
genes every member carries; Hoepfner covers 3,453 of the 3,587 background genes, which is
why the best-available core set is smaller. The composed top 100 is 100 by
construction.}
\label{tab:nominations}
%% GENERATED by experiments/040-inhibitor-synergy-wetlab/scripts/notes_tex_tables.py
%% SOURCE: the result files named in that script's docstring. Do not edit by hand.
\begin{tabular}{p{62mm}rlrl}
\toprule
Category & Median genes, & Range, & Median genes, & Range, \\
 & best available (n) & best available (n) & all predicted (n) & all predicted (n) \\
\midrule
Core, sensitive in at least four of six singles & 67 & 60 to 68 & 105 & 105 to 105 \\
Compound specific, sensitive in exactly one member & 65 & 7 to 125 & 22 & 2 to 64 \\
Conflict, sensitive in one member and resistant in another & 11 & 5 to 25 & 1 & 1 to 1 \\
Composed top 100 & 100 & 100 to 100 & 100 & 100 to 100 \\
\bottomrule
\end{tabular}
%% over 35 combinations, rule sum

\end{table}
